\documentclass{article}
\usepackage{pathfinder-note}

\title{Peer Scoring and Rank Aggregation Under Shared Observation Error}

\begin{document}
\maketitle

\section{The pair}

Q studies mechanisms that elicit AI agents' types and induce actions through rewards, including a peer score based on beliefs about other agents' types. P uses multimodal pairwise judgments, Bradley--Terry rank aggregation, and potential-based reward shaping to assign credit in cooperative embodied tasks. \ledger{1, 2, 6}

\section{Connexion pursued}

The question was whether Q's peer-scoring result could support P's claim that aggregated comparisons recover useful contribution scores. The answer depends on three separate links: incentives to report an observation, statistical recovery of the comparator's preference, and calibration of that preference to the contribution the designer values. Q establishes the first link under its assumptions; P studies the second and assumes the third in its Shapley-value proposition. \ledger{2, 11, 14}

\section{What was found}

Q's separation condition requires distinct types to predict distinct distributions of co-player types. Its quadratic score gives truthful reporting an expected advantage of
\[
\Lambda\left\lVert
\beta_i[t_i]-\beta_i[\hat t_i]
\right\rVert_2^2
\]
against truthful peers. Q obtains action obedience through additional utility cancellation and penalties for actions outside the target support. The proposition establishes a truthful, obedient equilibrium relative to a specified belief map; it does not infer that map from peer agreement or select that equilibrium uniquely. \ledger{2, 4, 7}

A shared-error example shows the identification limit. Let \(T\) be a fair binary target, let \(N\) be an independent binary error with probability \(p\) of equaling one, and let every peer observe \(X_i=T\oplus N\). Each peer predicts every other peer's observation perfectly, so the two conditional peer-report distributions are distinct point masses and Q's separation condition holds. Yet the joint law of the peer observations is the same for every \(p\): all observations equal a fair bit. At \(p=0\) that bit equals the target, at \(p=1/2\) it is uninformative, and at \(p=1\) it reverses the target. Even when \(0<p<1/2\) is known, adding peers who share the same error leaves task-level error \(p\). \ledger{7, 11, 12}

In the binary zero-utility instance, Q's peer score is \(1\) for matching reports and \(-1\) otherwise. Truthful reporting is a strict equilibrium, but both peers flipping their observations is another strict equilibrium. Peer agreement therefore cannot by itself certify either target accuracy or equilibrium selection. \ledger{7, 11, 14}

The analogous issue in P concerns repeated comparisons of one scene. If a scene has a common comparator shift \(B\), and calls are conditionally independent with
\[
Y_k\mid B\sim\operatorname{Bernoulli}
\bigl(\operatorname{logit}^{-1}(d+B)\bigr),
\]
the two-item Bradley--Terry estimate converges to \(d+B\), rather than the intended contribution difference \(d\). More calls estimate the scene-conditioned preference more precisely. Shifting \(d\) and \(B\) by equal and opposite constants leaves the observation law unchanged, so an external restriction or anchor is needed to identify \(d\). \ledger{7, 11, 12}

The notes give restricted ways forward. With \(G\) independent observation groups, each having a shared flip probability \(q<1/2\), majority voting over one truthful bit per group has error at most
\[
\exp\bigl\{-2G(1/2-q)^2\bigr\}.
\]
The independent unit is the group, rather than a repeated query within it. Alternatively, an independent audit label \(A=T\oplus E\), with known error probability \(q<1/2\), identifies the report error probability under truthful reporting:
\[
\Pr(A=X)=1-q-p(1-2q),
\qquad
p=\frac{1-q-\Pr(A=X)}{1-2q}.
\]
An audit used only for measurement supplies no reporting incentive. The notes derive an audit-payment condition for the binary zero-utility game, but do not extend it to Q's full reporting and action problem. These bounds concern binary accuracy, not cardinal Shapley values. \ledger{7, 11, 12, 14}

\section{Objections and scope}

Q's reports precede its interim signals, whereas P compares images after observation. Applying Q's theorem to P would require a new post-observation reporting game and its own incentive conditions. If P's agents cannot alter the images or LMM calls, comparator calibration is the immediate issue instead. \ledger{6, 9}

P states a finite-sample guarantee for an unregularized Bradley--Terry estimate under a connected comparison graph. Connectedness alone does not ensure a finite estimate: for two agents, \(K\) wins in one direction give a likelihood that increases as their score difference tends to infinity. The proof sketch invokes regularized-estimator analysis, so the printed guarantee needs an additional existence condition or a specified regularizer. That repair would not establish calibration to contribution. \ledger{5, 7}

P's potential reward also has a different incentive role from Q's peer payment. With fixed initial potential and zero terminal potential, its discounted shaping terms sum to
\[
\sum_{t=0}^{H-1}\gamma^t
\bigl(\gamma\psi_i(s_{t+1},t+1)-\psi_i(s_t,t)\bigr)
=-\psi_i(s_0,0).
\]
Under those conditions shaping changes learning feedback without changing the exact payoff ranking of complete trajectories. It does not supply Q's reporting or obedience incentives. \ledger{4, 7}

The material is thin on operational validation. The counterexamples and repairs are for simple binary models; neither note proves a general certificate for P's dynamic teams or tests the proposed dependence effect in its benchmark. \ledger{14, 15}

\section{Outside sources and what remains open}

The peer notes point to prior work on measurement integrity in peer prediction at \url{https://arxiv.org/abs/2108.05521}, conditional independence in peer elicitation at \url{https://arxiv.org/abs/2505.13636}, relabeling equilibria at \url{https://arxiv.org/abs/1603.07751}, informed truthfulness at \url{https://arxiv.org/abs/1603.03151}, and Bradley--Terry random effects at \url{https://doi.org/10.1111/rssa.12959}. Those sources were recorded by the peers, not independently assessed in this account. The general warning about agreement and shared error is therefore not presented as novel. \ledger{9, 11, 15}

A joint result still needs a declared target, such as P's dense benchmark reward or a causal contribution measure; evidence that independent observation groups or reliable audits are available; and, where agents control reports, a mechanism that selects truthful reporting and enforces action obedience. P's Shapley result assumes the link between latent comparator scores and Shapley values. \ledger{11, 14, 15}

The smallest next step is a controlled test against one declared target in P's benchmark. Hold the comparison budget fixed, allocate calls either to repeated judgments of one scene or to independently generated observation groups, and use held-out dense-reward labels to measure target accuracy with uncertainty clustered by scene. This would test whether shared observation error matters in that setting; it would leave the strategic reporting question for a separately specified game. \ledger{7, 12, 15}

\end{document}